\begin{chineseabstract}
脑机接口（BCI）物联网链路同时承载高敏感神经信号与控制指令，面临重放、篡改、中间人攻击与流量侧写等复合风险。针对受限设备、公网可达性和外观可识别性共同带来的协议设计难点，本文实现并验证了一种面向 BCI-IoT 的复合安全通信协议。

方法上，本文在应用层与 TCP 之间构建协议增强层：握手阶段采用证书签名、临时密钥协商、摘要绑定与确认码；会话阶段采用 AEAD 记录封装；外观层采用可逆字节映射与随机填充，并在纯模式和 6bit 分包模式之间形成可配置取舍。为保证论证严谨性，本文设置本地对照、侧写与主动攻击、跨区域受限网络三类实验，并在统一指标定义下开展多轮重复统计。

实验结果显示，所提方案在受限公网条件下可保持与主流安全协议同量级的时延，并在数据面内存占用上具有优势；在安全验证中可稳定拦截重放与篡改输入；在侧写对抗中，随机填充可显著降低长度序列侧写的类别恢复率。

本文工作给出了面向 BCI-IoT 的通信层协议原型及其复现实验框架。结果表明，随机填充能将被动侧写攻击的类别恢复率从 100\% 降至 16.67\%；代价是更高的字节流开销，致使带宽利用率降低。研究结论可为 BCI-IoT 安全通信协议的工程选型与后续标准化提供参考。

\chinesekeyword{脑机接口，物联网安全通信，数独编码，随机填充，流量侧写}
\end{chineseabstract}
